\centerline{\bf \Huge Перестановки}

\begin{flushright}
        \scriptsize{листочек для фарма баллов}
\end{flushright}

\textbf{Определение.} Пусть дано конечное множество $M$. (Часто $M=\{1,2, \ldots, n\}$ ). \textit{Перестановка} --- биективное отображение $\sigma: M \rightarrow M$.

Пример.

$$
\sigma=\left(\begin{array}{llllllllllll}
1 & 2 & 3 & 4 & 5 & 6 & 7 & 8 & 9 & 10 & 11 & 12 \\
3 & 7 & 2 & 6 & 1 & 9 & 5 & 8 & 4 & 10 & 12 & 11
\end{array}\right)
$$


Здесь, например, $\sigma(3)=2$, а $\sigma(8)=8$.

\textbf{Определение.} \textit{Инверсией} перестановки $\sigma$ называется пара чисел $i, j$ такая, что $i<j$, но $\sigma(i)>\sigma(j)$. Так, в примере пара $(2,4)$ образует инверсию, так как $2<4$, но $7>6$.

\textbf{Определение.} \textit{Четность} перестановки --- это четность числа ее инверсий.

\textbf{Определение.} \textit{Цикл} $\left(i_1, i_2, \ldots, i_k\right)$ - такая перестановка, что $\sigma\left(i_j\right)=i_{j+1} ; \sigma\left(i_k\right)=i_1$; $\sigma(l)=l$ при $l \notin\left\{i_1, i_2, \ldots, i_k\right\}$.

\textbf{Определение.} \textit{Транспозиция} --- цикл длины два.

Произвольную перестановку бывает удобно представить в виде композиции непересекающихся циклов. Так, перестановка из примера - композиция циклов ( $1,3,2,7,5$ ), $(4,6,9),(11,12),(8),(10)$. (В таком разложении принято выделять неподвижные элементы как циклы длины 1).

\textbf{Определение.} \textit{Тип} перестановки - набор длин циклов в ее разложении. Так, перестановка из примера имеет тип $(5,3,2,1,1)$.

\problem{1}
(a) Известны четности перестановок $\sigma$ и $\tau$. Что можно сказать о четности композиции $\sigma \circ \tau$ ?

\noindent
(b) Найдите четность перестановки типа ( $d_1, d_2, \ldots, d_k$ ).


\newpage

\hspace{3mm}

\problem{2}
(a) Докажите, что любую перестановку можно представить в виде композиции нескольких транспозиций.

\noindent
(b) Опишите все перестановки, которые можно представить в виде композиции циклов длины 3 .

\problem{3}
(a) Сколько существует перестановок на 10 -элементном множестве?

\noindent
(b) Каких перестановок среди них больше: четных или нечетных?

\noindent
(c) Сколько существует перестановок типа ( $1,2,3,4$ )?

\problem{4}\noindent
На полке в перепутанном порядке стоят $n$ пронумерованных томов собрания сочинений Владимира Ленина. Каждую минуту библиотекарь меняет местами (a) некоторые две соседние книги; (b) две произвольные книги.

За какое наименьшее время библиотекарь может гарантированно расставить книги в порядке возрастания?

\problem{5}
Нескольким детям дали по карандашу одного из трех цветов. Дети как-то поменялись карандашами, после чего у каждого оказался не тот карандаш, который был у него вначале. Докажите, что цвета карандашей могли быть такими, что у каждого вначале и в конце карандаши были разных цветов.

\problem{6}
В некотором городе разрешаются только парные обмены квартир (если две семьи обмениваются квартирами, то в тот же день они не имеют права участвовать в другом обмене). Докажите, что любой сложный обмен квартирами можно осуществить за два дня.

\problem{7}
Придумайте две перестановки, такие, что их композициями можно получить любую перестановку.

\problem{8}
После усердных тренировок Эльзята в совершенстве овладела навыком сборки кубика Рубика. Она уверяет, что может даже собрать кубик Рубика, у которого перевернули ровно один из боковых кубиков (то есть тот, что с двумя цветными гранями). Можно ли верить Эльзяте?